\section*{ID6812: From Linear Modes to Nonlinear Filaments: Zel’dovich Approximation: λ\_1,λ\_2}
\paragraph{Metadata.}
PiID: 3543947822181 | RTEID: RTE:P26452.5088:T5.8865 | UUID: e6dc051d-cd62-56c0-a1a7-8da73e275e11

\paragraph{Canonical Statement.}
So **mathematically**: - Compute \textbackslash{}(\textbackslash{}Phi(\textbackslash{}mathbf\{q\}) = \textbackslash{}sum\_\{i\} B\_i\textbackslash{}, j\_\{\textbackslash{}ell\_i\}(k\_i q)\textbackslash{}, Y\_\{\textbackslash{}ell\_i\}\textasciicircum{}\{m\_i\}(\textbackslash{}hat q)\textbackslash{}). - Compute \textbackslash{}(T\_\{ij\}=\textbackslash{}partial\_i\textbackslash{}partial\_j\textbackslash{}Phi\textbackslash{}) (use spherical-coordinate derivative formulas or transform to Cartesian with radial/angular chain rule). - Find eigenvalues \textbackslash{}(\textbackslash{}lambda\_1\textbackslash{}ge\textbackslash{}lambda\_2\textbackslash{}ge\textbackslash{}lambda\_3\textbackslash{}). Regions where \textbackslash{}(\textbackslash{}lambda\_1,\textbackslash{}lambda\_2\textbackslash{}) large and positive are filament seeds.

\paragraph{Canonical Equation.}
\[
\lambda_1,\lambda_2
\]

\paragraph{Dictionary.}
Relation: λ\_1,λ\_2

\paragraph{Encyclopedia.}
From Linear Modes to Nonlinear Filaments: Zel’dovich Approximation: λ\_1,λ\_2 is a symbolic expression, written λ\_1,λ\_2. It introduces a symbol or relation used elsewhere in the framework. Within the Trawin Zero Point registry it serves as one element of the framework's mathematical narrative.
